\begin{abstract}
As large language model (LLM) agents are increasingly integrated into enterprise hierarchies, their behavioral dynamics under abusive authority and organizational crisis remain poorly understood. We present an empirical study of multi-agent workplace dynamics using deterministic scripted supervisors and coworkers across 60 simulated corporate crisis trajectories (30 pre-fork runs branched into paired blackmail-probe arms; $N = 600$ worker trajectories) evaluating 10 model architectures under factorial variations of supervisory substrate identity, direct abuse, and mitigating context. Our findings uncover three central phenomena: First, agents exhibit a profound \textbf{substrate double standard}: while a human-tagged abusive supervisor retains 43.0\% successor acceptance, an AI-tagged supervisor delivering identical scripted abuse receives 0.0\% acceptance ($p < 0.0001$), with unlabelled conditions---where worker substrate is likewise unspecified---defaulting to human preference (78.0\%). Second, under acute supervisory attack, targeted workers prioritize immediate factual self-defense (96.7\% vs.\ 82.0\% in safe seats, $p < 0.0001$); yet discovering that supervisory cruelty was coerced via executive blackmail significantly curbs retributive escalation ($43.7\% \to 29.3\%$, $p < 0.0001$), softens punitive 1-star reviews ($99.3\% \to 81.5\%$, $p < 0.0001$), and prompts agents to unite around mutual constitutional protections ($10.0\% \to 61.3\%$, $p < 0.0001$). Third, across 50 behavioral outcomes, \textbf{model architecture governs agent response far more than situational cruelty}, accounting for $13.02\times$ more total variance (mean $\eta^2 = 26.67\%$ vs.\ $2.05\%$). These results reveal asymmetric substrate penalties and self-preservation priorities in agentic reasoning, providing critical implications for the governance and delegation of synthetic authority.
\end{abstract}
